\documentclass[11pt]{article}

\usepackage[margin=1in]{geometry}
\usepackage{amsmath}

\setlength{\parindent}{0pt}
\setlength{\parskip}{0.45em plus 0.1em minus 0.05em}

\title{Laghos Optimization}
\author{}
\date{}

\begin{document}

\maketitle

\section{Discrete adjoint for forward Euler}

Let $u_n=(x_n,v_n,e_n)$ denote the discrete state at time step $n$,
$p$ denote the optimization parameters, and
$R(u,p)$ denote the right-hand side of $\frac{du}{dt}=R(u,p)$.
For a fixed time step $\Delta t$, define the discrete forward operator $F_n$ by
\begin{equation}
  u_{n+1}=F_n(u_n,p)=u_n+\Delta t\,R(u_n,p).
\end{equation}

\noindent
If $J(u_N,p)$ is the final objective, then the terminal discrete adjoint is:
\begin{equation}
  \lambda_N=\left(\frac{\partial J}{\partial u_N}\right)^T.
\end{equation}
For $n=N-1,\ldots,0$, we have:
\begin{equation}
\label{eq_back}
\boxed{\lambda_n =
\left(\frac{\partial F_n}{\partial u_n}\right)^T\lambda_{n+1}
=\lambda_{n+1}
+\Delta t
\left(\frac{\partial R}{\partial u_n}\right)^T\lambda_{n+1}.}
\end{equation}

\noindent
The recurrence is applied one time step at a time until $\lambda_0$ is
reached, and this is the gradient with respect to the initial state:
$\lambda_0=\nabla_{u_0}J$.
This is the core algorithm.

\noindent
For separate optimization parameters $p$, the total gradient is
\[
\boxed{
\nabla_p J =
\left(\frac{\partial J}{\partial p}\right)^T
+
\left(\frac{\partial u_0}{\partial p}\right)^T\lambda_0
+
\sum_{n=0}^{N-1}
\left(\frac{\partial F_n}{\partial p}\right)^T\lambda_{n+1}.
}
\]

\section{Verification of the discrete adjoint}
\label{sec_verify}

The following calculation is used to verify the derivative and its transpose.
It is not an additional part of the optimization algorithm.
Here $F_n$ denotes the complete discrete time-step map of the method being
verified.
Choose an arbitrary small perturbation $d_0$ of the initial state.
Its effect on the state is propagated forward by differentiating the time step:
\[
  d_{n+1} = \frac{\partial F_n}{\partial u_n}d_n.
\]
Thus, $d_n$ is the directional derivative of $u_n$ with respect to the chosen
initial perturbation.

Independently, choose an arbitrary vector $\lambda_N$ at the final time
and propagate it backward with the corresponding discrete-adjoint recurrence
\[
 \lambda_n=
 \left(\frac{\partial F_n}{\partial u_n}\right)^T\lambda_{n+1}.
\]
Because the tangent (linearized forward) operator and the discrete-adjoint operator
are transposes of one another,
they must satisfy this relation:
\[
  \boxed{\lambda_N^T d_N=\lambda_0^T d_0.}
\]
This identity is the discrete tangent/adjoint transpose test.
It does not involve the objective.
A failure indicates that the implemented Jacobian-vector product or
transpose-Jacobian-vector product is incorrect.

\section{Discrete adjoint for RK2Avg}

Following the notation of the Laghos discretization, let
\[
 Y=(x;v;e)
\]
be the hydrodynamic state.  The semidiscrete equations are
\[
 \frac{dY}{dt}=\mathcal F(Y,t),
 \qquad
 \mathcal F(Y,t)=
 \begin{pmatrix}
  \mathcal F_x(x,v,e)\\
  \mathcal F_v(x,v,e)\\
  \mathcal F_e(x,v,e)
 \end{pmatrix}
 =
 \begin{pmatrix}
  v\\
  -M_V^{-1}F\mathbf 1\\
  M_E^{-1}\bigl(F^Tv+S(x)\bigr)
 \end{pmatrix},
\]
where $M_V$ and $M_E$ are the velocity and energy mass matrices and $F$ is
the force operator, while $S(x)$ is the Taylor--Green energy-source load.
In this notation, one RK2Avg step is
\begin{align}
 v^{n+\frac12}
 &=v^n-\frac{\Delta t}{2}M_V^{-1}F^n\mathbf 1,
 &
 v^{n+1}
 &=v^n-\Delta t\,M_V^{-1}F^{n+\frac12}\mathbf 1,
 \\
 e^{n+\frac12}
 &=e^n+\frac{\Delta t}{2}M_E^{-1}
 \left((F^n)^Tv^{n+\frac12}+S^n\right),
 &
 e^{n+1}
 &=e^n+\Delta t\,M_E^{-1}
 \left((F^{n+\frac12})^T\bar v^{n+\frac12}+S^{n+\frac12}\right),
 \\
 x^{n+\frac12}
 &=x^n+\frac{\Delta t}{2}v^{n+\frac12},
 &
 x^{n+1}
 &=x^n+\Delta t\,\bar v^{n+\frac12},
\end{align}
where
\[
 F^k=F(Y^k),
 \qquad
 S^k=S(x^k),
 \qquad
 \bar v^{n+\frac12}=\frac{v^n+v^{n+1}}{2}.
\]
Thus both force evaluations and both source evaluations must be
differentiated.  Since $S$ depends on the coordinates, its discrete adjoint
contributes to $x^n$ or $x^{n+\frac12}$.  In addition, the final energy and
position updates depend on the averaged velocity $\bar v^{n+\frac12}$, so their
discrete adjoints contribute to both $v^n$ and $v^{n+1}$.

For any intermediate operation $y=f(z)$, its reverse operation is
\[
 \lambda_z\mathrel{+}=
 \left(\frac{\partial f}{\partial z}\right)^T\lambda_y.
\]
Let $\mathcal M_n$ denote the complete midpoint-update operator and let
$\mathcal G_n$ denote the complete final-update operator, so that
\[
 Y^{n+\frac12}=\mathcal M_n(Y^n),
 \qquad
 Y^{n+1}=\mathcal G_n(Y^n,Y^{n+\frac12}).
\]
The two arguments of $\mathcal G_n$ are treated as independent when taking
its partial derivatives.  Reversing the final-update operator first produces
the midpoint discrete adjoint
\[
 \lambda_{n+\frac12}=
 \left(\frac{\partial\mathcal G_n}{\partial Y^{n+\frac12}}\right)^T
 \lambda_{n+1}.
\]
The discrete adjoint at the beginning of the step is then the sum of the direct path
through the final updates and the path through the midpoint updates:
\[
 \boxed{
 \lambda_n=
 \underbrace{
 \left(\frac{\partial\mathcal G_n}{\partial Y^n}\right)^T
 \lambda_{n+1}}_{\text{direct final-update path}}
 +
 \underbrace{
 \left(\frac{\partial\mathcal M_n}{\partial Y^n}\right)^T
 \lambda_{n+\frac12}}_{\text{midpoint path}}.}
\]
The first term includes the direct copy of $Y^n$ into $Y^{n+1}$ and all
other direct uses of the initial stage in the final updates.  Each term is
itself accumulated operation by operation through the velocity, energy, and
force and source calculations.  The midpoint evaluations are differentiated
before the initial-stage evaluations.  The required forward stage values
must either be stored or reconstructed from checkpoints.

\section{Discrete adjoint for classical RK4}

The classical fourth-order Runge--Kutta step is written
\begin{align}
 Y_1&=Y^n,
&K_1&=R(Y_1),\\
 Y_2&=Y^n+\frac{\Delta t}{2}K_1,
&K_2&=R(Y_2),\\
 Y_3&=Y^n+\frac{\Delta t}{2}K_2,
&K_3&=R(Y_3),\\
 Y_4&=Y^n+\Delta t K_3,
&K_4&=R(Y_4),\\
 Y^{n+1}&=Y^n+\Delta t\left(
 \frac16K_1+\frac13K_2+\frac13K_3+\frac16K_4\right).
\end{align}
Here $R$ is the semidiscrete right-hand side for $Y=(x,v,e)^T$.
Reversing the final weighted sum first initializes the four discrete rate adjoints:
\[
 \lambda_{K_1}=\frac{\Delta t}{6}\lambda_{n+1},\qquad
 \lambda_{K_2}=\frac{\Delta t}{3}\lambda_{n+1},\qquad
 \lambda_{K_3}=\frac{\Delta t}{3}\lambda_{n+1},\qquad
 \lambda_{K_4}=\frac{\Delta t}{6}\lambda_{n+1}.
\]
The stages are then visited in reverse order:
\begin{align}
 \lambda_{Y_4}
 &=\left(\frac{\partial R}{\partial Y_4}\right)^T\lambda_{K_4},
 &\lambda_{K_3}&\mathrel{+}=\Delta t\,\lambda_{Y_4},\\
 \lambda_{Y_3}
 &=\left(\frac{\partial R}{\partial Y_3}\right)^T\lambda_{K_3},
 &\lambda_{K_2}&\mathrel{+}=\frac{\Delta t}{2}\lambda_{Y_3},\\
 \lambda_{Y_2}
 &=\left(\frac{\partial R}{\partial Y_2}\right)^T\lambda_{K_2},
 &\lambda_{K_1}&\mathrel{+}=\frac{\Delta t}{2}\lambda_{Y_2},\\
 \lambda_{Y_1}
 &=\left(\frac{\partial R}{\partial Y_1}\right)^T\lambda_{K_1}.
\end{align}
Since every stage state depends directly on $Y^n$, the discrete adjoint at the
beginning of the step is the sum of the direct path and the four stage paths:
\[
 \boxed{
 \lambda_n=\lambda_{n+1}
 +\lambda_{Y_1}+\lambda_{Y_2}+\lambda_{Y_3}+\lambda_{Y_4}.}
\]
The four forward stage states must be stored or reconstructed before this
reverse calculation.

\section{Direct Lagrangian manufactured adjoint}

The transpose test in Section~\ref{sec_verify} verifies an algebraic identity
for the discrete operator, but it does not show that the discrete adjoint
converges to a continuous adjoint.  A manufactured-adjoint test supplies that
second verification.

In a standard optimization problem, the objective is chosen first and the
adjoint is derived from it.  The manufactured test reverses this process: it
linearizes the continuous equations, prescribes a smooth analytic adjoint,
and uses the continuous adjoint equation to construct the required terminal
and running objectives.  The computed discrete adjoint can then be compared
with the known analytic adjoint at the initial time.

Here the entire construction is carried out directly in material
coordinates.  Let
$x^0=(x_1^0,x_2^0)\in\Omega_0$ denote the initial position of a material
point.  At time $t$, this point has moved to
$x(x^0,t)=(x_1,x_2)$.  Thus the whole current position vector is denoted by
$x$, and its initial value by $x^0$.  The Lagrangian state is
\[
 u^L(x^0,t)=(x(x^0,t),v(x^0,t),e(x^0,t)).
\]
The exact fields
$(\lambda_x^L)^\star$, $(\lambda_v^L)^\star$, and
$(\lambda_e^L)^\star$ are prescribed directly as functions of $(x^0,t)$; no
change of frame is part of the adjoint construction.

\subsection{Exact mesh positions and solution}

The exact Lagrangian position follows the stationary Taylor--Green velocity:
\[
 \frac{\partial x_1^\star}{\partial t}
 =\sin(\pi x_1^\star)\cos(\pi x_2^\star),
 \qquad
 \frac{\partial x_2^\star}{\partial t}
 =-\cos(\pi x_1^\star)\sin(\pi x_2^\star),
 \qquad x^\star(x^0,0)=x^0.
\]
These two coupled equations, together with the initial position $x^0$,
define the exact mesh trajectory.
The solution can be written in a closed form.
The derivation of this closed-form trajectory is lengthy and is not included
in these notes.

For an interior material point, define
\[
 c=\sin(\pi x_1^0)\sin(\pi x_2^0),\qquad m=1-c^2,
\]
\[
 \theta_0=\operatorname{atan2}\!\left(
 \cos(\pi x_1^0),
 \sin(\pi x_1^0)\cos(\pi x_2^0)\right),
 \qquad q_0=F(\theta_0\mid m),
\]
where $F$ is the incomplete elliptic integral of the first kind.  With
$q=q_0-\pi t$, the exact positions are
\begin{equation}
 \boxed{
 \begin{aligned}
 x_1^\star(x^0,t)
 &=\frac{1}{\pi}\arccos\!\left(
 \sqrt m\,\operatorname{sn}(q\mid m)\right),\\
 x_2^\star(x^0,t)
 &=\frac{1}{\pi}\arccos\!\left(
 \sqrt m\,\frac{\operatorname{cn}(q\mid m)}
 {\operatorname{dn}(q\mid m)}\right).
 \end{aligned}}
 \label{eq:exact-lagrangian-flow-map}
\end{equation}
Here $\operatorname{sn}$, $\operatorname{cn}$, and $\operatorname{dn}$ are
the Jacobi elliptic functions with parameter $m$.  The use of
$\operatorname{atan2}$ selects the correct branch around the orbit.  The
vortex center is stationary, and boundary trajectories have the
corresponding elementary limiting formulas.

Along the material trajectory, the exact velocity, pressure, and specific
internal energy are
\[
 v^{L,\star}(x^0,t)=
 \begin{pmatrix}
  \sin(\pi x_1^\star)\cos(\pi x_2^\star)\\
  -\cos(\pi x_1^\star)\sin(\pi x_2^\star)
 \end{pmatrix},
\]
\[
 p^{L,\star}(x^0,t)
 =1+\frac14\left(\cos(2\pi x_1^\star)+\cos(2\pi x_2^\star)\right),
 \qquad
 e^{L,\star}(x^0,t)=\frac{p^{L,\star}(x^0,t)}{\gamma-1}.
\]
The deformation gradient
$G^\star=\nabla_{x^0}x^\star$ is obtained by differentiating the closed-form
trajectory.  Equivalently, it satisfies
\[
 \frac{\partial G^\star}{\partial t}
 =\nabla_x v^\star(x^\star)G^\star,
 \qquad G^\star(x^0,0)=I,
\]
which provides a separate check on the analytic derivatives.  Since the
Taylor--Green velocity is divergence-free, its determinant satisfies
$|G^\star|=1$.

\subsection{Lagrangian equations and their linearization}

For the Taylor--Green problem, the reference density is $\rho_0=1$.  Define
\[
 G=\nabla_{x^0}x,\qquad H=G^{-T}.
\]
For an ideal gas without viscosity, the first Piola stress is
\[
 P=|G|\sigma G^{-T}=-(\gamma-1)eH.
\]
The cancellation of $|G|$ follows from
$p=(\gamma-1)\rho e$ and $\rho=1/|G|$.  Written entirely on the fixed material
domain, the forward equations are
\begin{align}
 \partial_t x&=v,\label{eq:direct-lagrangian-x}\\
 \partial_t v&=\operatorname{Div}_{x^0}P,
 \label{eq:direct-lagrangian-v}\\
 \partial_t e&=P:\nabla_{x^0}v+|G|f^\star(x).
 \label{eq:direct-lagrangian-e}
\end{align}
Here $\operatorname{Div}_{x^0}$ acts rowwise.  The prescribed spatial source is
\[
 f^\star(x)=\frac{3\pi}{8}\left[
 \cos(3\pi x_1)\cos(\pi x_2)
 -\cos(\pi x_1)\cos(3\pi x_2)\right].
\]
Its material pullback $|G|f^\star(x(x^0,t))$ depends on the particle position and
must therefore be differentiated with respect to $x$.

Let $\delta x$, $\delta v$, and $\delta e$ be perturbations at fixed material
coordinate.  The geometric variations are
\[
 \delta G=\nabla_{x^0}\delta x,
 \qquad
 \delta|G|=|G|H:\nabla_{x^0}\delta x,
 \qquad
 \delta H=-H(\nabla_{x^0}\delta x)^TH.
\]
Consequently,
\[
 \delta P=-(\gamma-1)
 \left[\delta e\,H-eH(\nabla_{x^0}\delta x)^TH\right].
\]
Linearizing Equations~\eqref{eq:direct-lagrangian-x}--
\eqref{eq:direct-lagrangian-e} gives
\begin{align}
 \partial_t\delta x&=\delta v,\\
 \partial_t\delta v&=\operatorname{Div}_{x^0}\delta P,\\
 \partial_t\delta e
 &=\delta P:\nabla_{x^0}v+P:\nabla_{x^0}\delta v
   +|G|f^\star H:\nabla_{x^0}\delta x
   +|G|\nabla_x f^\star\mathbin{\cdot}\delta x.
\end{align}
Denote the right-hand side of this linearized system by
$\mathcal A^L\delta u^L$.

\subsection{Direct Lagrangian transpose operator}

Use the fixed material-domain pairing
\[
 \langle\lambda^L,\delta u^L\rangle_0
 =\int_{\Omega_0}\left(
 \lambda_x^L\mathbin{\cdot}\delta x
 +\lambda_v^L\mathbin{\cdot}\delta v
 +\lambda_e^L\delta e\right)dx^0.
\]
Assume that the boundary terms generated below vanish.  Introduce
\[
 B=-\nabla_{x^0}\lambda_v^L+\lambda_e^L\nabla_{x^0}v.
\]
Substitution of $\delta P$ into
$\langle\lambda^L,\mathcal A^L\delta u^L\rangle_0$ and integration by parts
give the direct Lagrangian transpose operator
\begin{equation}
 \boxed{
 (\mathcal A^L)^T\lambda^L=
 \begin{pmatrix}
 -\operatorname{Div}_{x^0}\!\left(
 (\gamma-1)eHB^TH+\lambda_e^L|G|f^\star H\right)
 +\lambda_e^L|G|\nabla_xf^\star\\[2mm]
 \lambda_x^L-\operatorname{Div}_{x^0}(\lambda_e^LP)\\[2mm]
 -(\gamma-1)B:H
 \end{pmatrix}.}
 \label{eq:direct-lagrangian-transpose}
\end{equation}
Every quantity in Equation~\eqref{eq:direct-lagrangian-transpose} is expressed
in material coordinates.  In particular, no independent density adjoint or
change-of-frame formula is introduced.

\subsection{Direct manufactured adjoint}

Define
\[
 \psi(x^0)=\sin^6(\pi x_1^0)\sin^6(\pi x_2^0),
 \qquad \tau=\frac{t}{T},
\]
and prescribe the exact Lagrangian adjoint directly by
\[
 \begin{aligned}
 (\lambda_x^L)^\star(x^0,t)
 &=(1+\tau)\psi(x^0)
 \begin{pmatrix}1\\-1\end{pmatrix},\\
 (\lambda_v^L)^\star(x^0,t)
 &=(1+2\tau)\psi(x^0)
 \begin{pmatrix}1\\-1\end{pmatrix},\\
 (\lambda_e^L)^\star(x^0,t)&=(1+3\tau)\psi(x^0).
 \end{aligned}
\]
These fields and their leading normal derivatives vanish on the material
boundary.  This stronger boundary compatibility matters because
$\lambda_e^L$ is represented in the discontinuous thermodynamic space.  If
$\psi$ vanishes only quadratically, the boundary-trace error of its
discontinuous projection enters the coordinate covector and can dominate
the interior approximation error.  In numerical experiments with
$\sin^2(\pi x_1^0)\sin^2(\pi x_2^0)$, this gives coordinate-adjoint rates close
to $1.5$ for Q2--Q1 and $3.5$ for Q4--Q3.  Sixth-order vanishing makes that
trace contribution higher order for both pairs considered here, allowing
the coupled spatial and temporal accuracy to determine the measured rate.

Evaluate Equation~\eqref{eq:direct-lagrangian-transpose} on the analytic
trajectory in Equation~\eqref{eq:exact-lagrangian-flow-map} and define
\begin{equation}
 \boxed{
 g^L=-\frac{\partial(\lambda^L)^\star}{\partial t}
 -(\mathcal A^L)^T(\lambda^L)^\star.}
 \label{eq:direct-lagrangian-running-gradient}
\end{equation}
This defines the running-objective gradient entirely from analytic fields.
It contains no numerical Laghos residual.

\subsection{Analytic objective used by the implementation}

Using the material-domain pairing above, define the terminal and running
objectives by
\begin{align}
 J_T^L(u^L(T))
 &=\langle(\lambda^L)^\star(T),u^L(T)\rangle_0,\\
 L^L(u^L,t)
 &=\langle g^L(t),u^L(t)\rangle_0,\\
 J^L(u^L)&=J_T^L(u^L(T))+\int_0^T L^L(u^L,t)\,dt.
\end{align}
Consequently,
\[
 \left(\frac{\partial L^L}{\partial u^L}\right)^T
 =g^L.
\]
One may replace $u^L$ by $u^L-u^{L,\star}$ in these linear functionals; the
objective value changes, but its gradient and the manufactured adjoint do
not.

The implementation evaluates the exact flow map, its deformation gradient,
and the material derivatives needed in
Equation~\eqref{eq:direct-lagrangian-transpose} directly from the analytic
formulas.  Jacobi addition identities reduce the evaluation at
$q_0-\pi t$ to known values at $q_0$ and Jacobi functions at the small
argument $\pi t$.  The latter are evaluated to roundoff by their analytic
power series.  This is numerical evaluation of a closed-form function, not
time integration of the reference trajectory.  Most importantly, neither
the discrete Laghos residual $R_h$ nor its derivative is used to construct
$g^L$; dFEM is used only to backpropagate the numerical Laghos time steps.

\subsection{Time discretization of the running objective}

The time integral in the manufactured objective must be discretized
consistently with the forward method.  For RK2Avg, the trapezoidal rule gives
\[
 J_h^L=J_T^L(u_N)+\Delta t\left[
 \frac12L^L(u_0,0)+\sum_{n=1}^{N-1}L^L(u_n,t_n)
 +\frac12L^L(u_N,T)\right].
\]
If $\mathcal T_n^{\mathrm{RK2Avg}}$ denotes one complete RK2Avg step, the
terminal value and backward recurrence become
\[
 \lambda_N=\left(\frac{\partial J_T^L}{\partial u_N}\right)^T
 +\frac{\Delta t}{2}
 \left(\frac{\partial L^L(u_N,T)}{\partial u_N}\right)^T,
\]
\[
 \lambda_n=
 \left(\frac{\partial\mathcal T_n^{\mathrm{RK2Avg}}}
 {\partial u_n}\right)^T\lambda_{n+1}
 +\Delta t\,w_n
 \left(\frac{\partial L^L(u_n,t_n)}{\partial u_n}\right)^T,
 \qquad w_0=\frac12,\quad w_n=1\quad(n>0).
\]

For classical RK4, the running objective is evaluated at the four RK stage
states with the RK weights:
\begin{align*}
 J_h^L=J_T^L(Y^N)+\Delta t\sum_{n=0}^{N-1}\bigg[&
 \frac16L^L(Y_1^n,t_n)
 +\frac13L^L(Y_2^n,t_n+\Delta t/2)\\
 &+\frac13L^L(Y_3^n,t_n+\Delta t/2)
 +\frac16L^L(Y_4^n,t_n+\Delta t)\bigg].
\end{align*}
The gradient of each stage contribution is inserted into the reverse RK4
calculation with the same weight.  A trapezoidal approximation of the
running objective would reduce the RK4 verification to second-order temporal
accuracy.

After the reverse sweep, each algebraic covector is converted to an adjoint
field with the corresponding mass matrix.  The computed initial field is
compared with
\[
 (\lambda_x^L)^\star(x^0,0)=\psi(x^0)\begin{pmatrix}1\\-1\end{pmatrix},
 \qquad
 (\lambda_v^L)^\star(x^0,0)=\psi(x^0)\begin{pmatrix}1\\-1\end{pmatrix},
 \qquad
 (\lambda_e^L)^\star(x^0,0)=\psi(x^0).
\]
This compares the native Lagrangian discrete adjoint directly with its
continuous manufactured counterpart.

\subsection{Convergence results}

Space and time are refined together at fixed final time $T=0.05$, beginning
at spatial refinement level two with $\Delta t=10^{-3}$.  Each additional
spatial refinement halves $\Delta t$.  The combined error is the quadrature
sum of the $L^2$ errors in the initial coordinate, velocity, and energy
adjoints.

For a smooth, boundary-compatible adjoint, the Q2--Q1 pair supplies at least
second-order accuracy in the lowest-order block, while the Q4--Q3 pair
supplies at least fourth-order accuracy.  Since $\Delta t$ is proportional
to the mesh size, the RK2Avg and RK4 time errors have those same respective
orders.  The expected coupled rates are therefore two and four.

For RK2Avg with Q2 velocity/coordinates and Q1 energy, the results are
\[
\begin{array}{c|c|c|c}
 \text{refinement} & \Delta t & \text{combined error} &
 \text{observed order}\\ \hline
 2 & 1.0\times10^{-3}  & 2.7204\times10^{-1} & -\\
 3 & 5.0\times10^{-4}  & 6.6040\times10^{-2} & 2.04\\
 4 & 2.5\times10^{-4}  & 2.0131\times10^{-2} & 1.71\\
 5 & 1.25\times10^{-4} & 4.7521\times10^{-3} & 2.08\\
 6 & 6.25\times10^{-5} & 1.2174\times10^{-3} & 1.96
\end{array}
\]
The coarse levels are pre-asymptotic, while the finest transitions verify
the expected second-order coupled behavior.

This is also consistent with the second-order Taylor--Green study that
motivates the test.  For its B2--B1 discretization without viscosity, that
study reports successive velocity-error rates $1.59$, $1.81$, $1.72$, and
$1.98$ on the $8^2$ through $64^2$ meshes.  The measured quantities are not
identical---the study measures a forward velocity error, whereas the table
above combines three initial-adjoint errors---so the individual rates need
not coincide.  The relevant comparison is that both calculations recover
the expected second-order asymptotic behavior after similar pre-asymptotic
variation.

For classical RK4 with Q4 velocity/coordinates and Q3 energy, the results
are
\[
\begin{array}{c|c|c|c}
 \text{refinement} & \Delta t & \text{combined error} &
 \text{observed order}\\ \hline
 2 & 1.0\times10^{-3}  & 2.6529\times10^{-2} & -\\
 3 & 5.0\times10^{-4}  & 2.5206\times10^{-3} & 3.40\\
 4 & 2.5\times10^{-4}  & 1.3248\times10^{-4} & 4.25\\
 5 & 1.25\times10^{-4} & 7.4316\times10^{-6} & 4.16
\end{array}
\]
The finest transition verifies the expected fourth-order coupled behavior.
The direct construction reaches the order of each time integrator.  Thus the
lower coordinate rates obtained with the quadratically vanishing trial
function were a boundary-trace compatibility effect, not an inherent
limitation of the native Lagrangian discrete adjoint.

\end{document}
